\section{Theoretical Setup}\label{sec:setup}

\subsection{Tensor model and objective}

Fix constants $\kappa\geq 1$ and $q>0$, and let $r,n,k$ be positive
integers.  For each mode, let
$\bar A=[\bar a_1,\ldots,\bar a_r]\in\mathbb R^{n\times r}$, with
analogous matrices $\bar B$ and $\bar C$.  For
$\bar M\in\{\bar A,\bar B,\bar C\}$, write
\[
D_{\bar M}=\operatorname{diag}(\|\bar m_1\|_2,\ldots,
\|\bar m_r\|_2),
\qquad
\widetilde M=\bar M D_{\bar M}^{-1}.
\]
Set $\rho=r^{-q}$.  Independently over $j\in[r]$ and the three modes,
draw
\[
\xi_j^a,\xi_j^b,\xi_j^c
\sim \mathcal N\!\left(0,\frac{\rho^2}{n}I_n\right),
\]
and define
\[
a_j=\bar a_j+\xi_j^a,
\qquad b_j=\bar b_j+\xi_j^b,
\qquad c_j=\bar c_j+\xi_j^c.
\]
The realized target tensor is
\[
T=\sum_{j=1}^r a_j\otimes b_j\otimes c_j
\in\mathbb R^{n\times n\times n}.
\]

For factor matrices $X=[x_1,\ldots,x_k]$, $Y=[y_1,\ldots,y_k]$, and
$Z=[z_1,\ldots,z_k]$, define
\[
\widehat T(X,Y,Z)=\sum_{i=1}^k x_i\otimes y_i\otimes z_i,
\qquad
\mathcal L(X,Y,Z)=\|T-\widehat T(X,Y,Z)\|_F^2.
\]
At time $t$, abbreviate
$\widehat T_t=\widehat T(X_t,Y_t,Z_t)$.  Fix one componentwise
Khatri--Rao ordering, put
\[
U_t^x=Z_t\odot Y_t,
\qquad U_t^y=Z_t\odot X_t,
\qquad U_t^z=Y_t\odot X_t,
\]
and let $T_{(1)},T_{(2)},T_{(3)}$ be the corresponding matricizations.

\subsection{Half-relaxed parallel ALS and product gauge}

All three least-squares candidates use the old iterate and are computed in
parallel:
\[
\begin{aligned}
X_{t+1}^{\mathrm{ls}}
 &=T_{(1)}U_t^x\big((U_t^x)^{\mathsf T}U_t^x\big)^\dagger,\\
Y_{t+1}^{\mathrm{ls}}
 &=T_{(2)}U_t^y\big((U_t^y)^{\mathsf T}U_t^y\big)^\dagger,\\
Z_{t+1}^{\mathrm{ls}}
 &=T_{(3)}U_t^z\big((U_t^z)^{\mathsf T}U_t^z\big)^\dagger.
\end{aligned}
\]
The Moore--Penrose pseudoinverse selects the minimum-Frobenius-norm
least-squares candidate.  With relaxation parameter $\eta=1/2$, form
\[
X_{t+1}^{\mathrm{raw}}=\tfrac12(X_t+X_{t+1}^{\mathrm{ls}}),
\quad
Y_{t+1}^{\mathrm{raw}}=\tfrac12(Y_t+Y_{t+1}^{\mathrm{ls}}),
\quad
Z_{t+1}^{\mathrm{raw}}=\tfrac12(Z_t+Z_{t+1}^{\mathrm{ls}}).
\]
There is no fixed-subspace constraint, clipping, regularizer, adaptive
preconditioner, restart, or early stopping.  Sequential Gauss--Seidel ALS
and unrelaxed simultaneous ALS are not covered.

After each relaxed sweep, apply the following deterministic gauge to every
component.  If $u=\|x_i\|_2$, $v=\|y_i\|_2$, and $w=\|z_i\|_2$ are
positive, let $g=(uvw)^{1/3}$ and replace
\[
(x_i,y_i,z_i)
\quad\text{by}\quad
(gx_i/u,gy_i/v,gz_i/w).
\]
The three scale factors have product one.  If $uvw=0$, replace the component
triple by $(0,0,0)$.  The resulting factors are denoted
$(X_{t+1},Y_{t+1},Z_{t+1})$.  The initial state is the raw Gaussian draw
specified below; the gauge is first applied after the first sweep.

\subsection{Primitive assumptions}

\begin{assumption}[Ambient dimension]\label{assump:dimension}
The integers $r,n$ are positive and
\[
n\geq C_{\mathrm{dim}}(\kappa,q)r^4\log r,
\]
where $C_{\mathrm{dim}}(\kappa,q)>0$ is independent of $r,n,k$ and of
the deterministic base triple.
\end{assumption}

\begin{assumption}[Full superlinear rank window]\label{assump:rank_window}
The integer $k$ satisfies
\[
r<k\leq r^{5/4}.
\]
Thus the universal overparameterization exponent is $1/4$.
\end{assumption}

\begin{assumption}[Well-conditioned deterministic bases]\label{assump:base_conditioning}
For every $\bar M\in\{\bar A,\bar B,\bar C\}$ and every $j\in[r]$,
\[
\kappa^{-1}\leq\|\bar m_j\|_2\leq\kappa,
\]
and all $r$ singular values of the rectangular matrix
$\widetilde M\in\mathbb R^{n\times r}$---equivalently, the square roots
of all $r$ eigenvalues of $\widetilde M^{\mathsf T}\widetilde M$---lie in
$[\kappa^{-1},\kappa]$.  This is the full-column-rank rectangular
singular-value convention used throughout.
\end{assumption}

\begin{assumption}[Independent Gaussian smoothing]\label{assump:gaussian_smoothing}
The exponent $q>0$ is fixed, $\rho=r^{-q}$, and the $3r$ perturbation
vectors have the mutually independent Gaussian law specified above.
\end{assumption}

\begin{assumption}[Independent Gaussian initialization]\label{assump:independent_initialization}
All $3nk$ entries of $X_0,Y_0,Z_0$ are mutually independent
$\mathcal N(0,1/n)$ variables and are independent of all smoothing
variables.
\end{assumption}

All probabilities below refer to this joint smoothing-and-initialization law,
conditional on the deterministic base triple.

\subsection{Formalized goal}

The goal is an explicitly conditional theorem.  We seek constants depending
only on $(\kappa,q)$ for which membership in the four-clause certificate
defined in Section~\ref{sec:preliminaries} implies that the objective has a
finite limit and that this limit is a positive constant fraction of
$\|T\|_F^2$.  The target fraction is
\[
\epsilon=
\left(\frac{\delta-L_P-\zeta}{\kappa^6C_T}\right)^2.
\]
The theorem does not assert that the certificate is nonempty, and it does not
assert any positive lower bound on the probability of the certificate.
Establishing such a uniform probability bound is the residual source-level
gap.
